\documentclass[11pt,a4paper]{article}
\usepackage[utf8]{inputenc}
\usepackage[T1]{fontenc}
\usepackage[brazil]{babel}
\usepackage{lmodern}
\usepackage[margin=2cm]{geometry}
\usepackage{amsmath,amssymb}
\usepackage{booktabs}
\usepackage{graphicx}
\usepackage{hyperref}
\usepackage{xcolor}
\usepackage{titlesec}
\usepackage{tabularx}
\usepackage{enumitem}
\usepackage{float}
\usepackage{array}
\renewcommand{\tabularxcolumn}[1]{>{\raggedright\arraybackslash}p{#1}}
\usepackage{newunicodechar}
\usepackage[most]{tcolorbox}
\usepackage{etoolbox}

\hypersetup{colorlinks=true, linkcolor=blue!70!black, citecolor=blue!70!black, urlcolor=blue!70!black}

\titleformat{\section}{\large\bfseries}{\thesection.}{0.5em}{}
\titleformat{\subsection}{\normalsize\bfseries}{\thesubsection.}{0.5em}{}
\titlespacing*{\section}{0pt}{1.2ex plus 0.5ex minus .2ex}{0.6ex}
\titlespacing*{\subsection}{0pt}{1.0ex plus 0.4ex minus .2ex}{0.4ex}
\setlength{\parskip}{0.45em}
\setlength{\parindent}{0pt}
\setlist{itemsep=0.15em,topsep=0.3em}

% --- símbolos ---
\newunicodechar{π}{\ensuremath{\pi}}
\newunicodechar{τ}{\ensuremath{\tau}}
\newunicodechar{φ}{\ensuremath{\varphi}}
\newunicodechar{ψ}{\ensuremath{\psi}}
\newunicodechar{Σ}{\ensuremath{\Sigma}}
\newunicodechar{→}{\ensuremath{\rightarrow}}
\newunicodechar{⇒}{\ensuremath{\Rightarrow}}
\newunicodechar{∈}{\ensuremath{\in}}
\newunicodechar{−}{\ensuremath{-}}
\newunicodechar{≈}{\ensuremath{\approx}}
\newunicodechar{≠}{\ensuremath{\neq}}
\newunicodechar{≡}{\ensuremath{\equiv}}
\newunicodechar{≤}{\ensuremath{\leq}}
\newunicodechar{≥}{\ensuremath{\geq}}
\newunicodechar{′}{\ensuremath{'}}
\newunicodechar{‑}{-}
\newunicodechar{✔}{\ensuremath{\checkmark}}
\newunicodechar{✘}{\ensuremath{\times}}

% --- cores e caixas ---
\definecolor{cC}{HTML}{7A4FA3}
\definecolor{cS}{HTML}{2F7D74}
\definecolor{cAcc}{HTML}{2F5D8A}
\newcommand{\metrictag}[2]{{\setlength{\fboxsep}{1.5pt}%
  \ifstrequal{#1}{c}{\colorbox{cC}{\textcolor{white}{\scriptsize\textbf{#2}}}}{%
  \ifstrequal{#1}{s}{\colorbox{cS}{\textcolor{white}{\scriptsize\textbf{#2}}}}{%
  \colorbox{black!60}{\textcolor{white}{\scriptsize\textbf{#2}}}}}}}
\newcommand{\pagechip}[1]{{\setlength{\fboxsep}{1.5pt}\colorbox{cAcc}{\textcolor{white}{\scriptsize\texttt{#1}}}}}
\newcommand{\readtag}[1]{\texttt{\small #1}}
\newcommand{\partbar}[1]{\par\vspace{1.2em}\noindent\colorbox{black!80}{\parbox{\dimexpr\linewidth-2\fboxsep}{\color{white}\bfseries\sffamily\MakeUppercase{#1}}}\par\vspace{0.4em}}
\newtcolorbox{citacard}[1]{breakable, enhanced, colback=white, colframe=black!25, boxrule=0.5pt, arc=2pt,
  left=6pt, right=6pt, top=4pt, bottom=4pt, title={\small #1}, coltitle=black, colbacktitle=black!6, fonttitle=\normalfont}
\newenvironment{origquote}{\par\itshape\small}{\par}
\newenvironment{tradquote}{\par\small\textsc{Tradução:}\ }{\par}
\newcommand{\explab}[1]{\par\smallskip{\small\color{cAcc}\textbf{\MakeUppercase{#1}}}\par}
\colorlet{notebg}{blue!4}\colorlet{notefg}{cAcc}
\colorlet{okbg}{green!5}\colorlet{okfg}{green!45!black}
\colorlet{warnbg}{orange!8}\colorlet{warnfg}{orange!80!black}
\newtcolorbox{infobox}[2][note]{breakable, enhanced, boxrule=0pt, arc=2pt, left=6pt, right=6pt, top=4pt, bottom=4pt,
  colback=#1bg, borderline west={3pt}{0pt}{#1fg}, frame hidden,
  before upper={\ifstrempty{#2}{}{\textbf{#2}\par}}}
\colorlet{pairc}{cC}\colorlet{pairs}{cS}
\newenvironment{paircol}[2]{\begin{tcolorbox}[width=0.485\linewidth, nobeforeafter, box align=top, boxrule=0.5pt,
  colframe=black!25, colback=white, borderline north={3pt}{0pt}{pair#1}, left=5pt, right=5pt,
  before upper={{\color{pair#1}\small\textbf{\MakeUppercase{#2}}}\par}]}{\end{tcolorbox}}

\newunicodechar{ℓ}{\ensuremath{\ell}}

\begin{document}

\begin{center}
    {\LARGE \textbf{Análise: a complexidade de Bhatt et al. (2020) no caso do ECG}} \\[0.35em]
    {\large BHATT, U.; WELLER, A.; MOURA, J. M. F. \textit{Evaluating and Aggregating Feature-based Model Explanations}. IJCAI, 2020, p.~3016--3022} \\[0.3em]
    \textbf{Métrica no Quantus:} \texttt{quantus.Complexity} \\[0.2em]
    \textbf{Complementa:} Fichamento da Faithfulness Correlation (\texttt{fichamentos/bhatt-2020-faithfulness-correlation}) e Seleção de métricas (\texttt{fichamentos/selecao-de-metricas}) \\[0.2em]
    \textbf{Aluno:} Arthur Santos de Oliveira (RA: 156297) \\[0.2em]
    \textbf{Atualizado em:} 23/10/2026 (primeira versão: 20/10/2026) \\[0.2em]
    \textbf{Repositório:} \url{https://github.com/Arthur-so/shap-ecg-arrhythmia}
\end{center}
\vspace{-0.5em}
\hrule
\vspace{0.8em}

\textbf{Por que este documento.} O fichamento da Faithfulness Correlation diz que o artigo propõe três critérios (baixa sensibilidade, alta fidelidade e baixa complexidade) e que “a complexidade não é usada”. Aqui se analisa se valeria a pena usá-la no TCC.

\textbf{O que mudou nesta versão (23/10).} Depois da primeira versão, a seleção de métricas (documento \textit{Seleção de métricas}, 22/10) escolheu as próximas métricas pela dimensão nova que trazem: Random Logit, eMPRT e esta complexidade. A localização clínica foi avaliada e deixada de fora. Esta versão registra a decisão, compara a métrica com as outras duas de complexidade do Quantus (seção~\ref{sec:escolha}), relaciona o resultado com o piso do GradientSHAP medido na sensibilidade (seção~\ref{sec:piso}) e com o eMPRT, que usa a mesma entropia (seção~\ref{sec:emprt}), e traz a implementação (\texttt{src/complexity.py}) com os resultados oficiais e a verificação pelo índice de Gini (seção~\ref{sec:impl}).

\textbf{Como ler.} \texttt{[p. N]} indica a página do PDF do artigo (1 a 7). As citações estão no original em inglês, cada uma seguida de uma tradução livre. A \textbf{Parte I} explica a métrica; a \textbf{Parte II} traz o que o artigo diz sobre ela; a \textbf{Parte III} mostra um diagnóstico com os mapas do modelo B, a escolha entre as métricas de complexidade e a decisão.

\begin{infobox}[ok]{Resposta curta}
\textbf{Incluída no TCC, como métrica complementar e na versão por regiões.} Ela mede algo que as métricas de fidelidade não medem (se a explicação é concentrada em poucas regiões, e portanto fácil de ler), custa quase nada (não chama o modelo) e completa os três critérios do mesmo artigo. A versão do Quantus, ponto a ponto, \textbf{não} serve: ela mede o comprimento do sinal, e não a explicação (seção~\ref{sec:pontos}). Entre as três métricas de complexidade do Quantus, a entropia de Bhatt et al.\ é a escolhida; o índice de Gini (Sparseness) entra só como verificação, e a Effective Complexity fica de fora (seção~\ref{sec:escolha}). A inclusão amplia o escopo do TCC para além de fidelidade e robustez e foi decidida junto com Random Logit e eMPRT; a mudança de rumo vai ao orientador.
\end{infobox}

\partbar{Parte I · Entendendo a métrica}

\section{O que a complexidade mede}
Um mapa SHAP dá uma importância a cada atributo da entrada. Se a importância está \textbf{espalhada} por quase todos os atributos, a explicação é difícil de ler: não há como apontar “o modelo olhou aqui”. Se está \textbf{concentrada} em poucos atributos, a explicação é simples. A complexidade mede esse grau de espalhamento.

\begin{infobox}[note]{Definição em uma frase}
Transformam-se as importâncias (em valor absoluto) em frações que somam 1, e calcula-se a \textbf{entropia} dessas frações. Entropia alta = importância espalhada = explicação complexa; entropia baixa = importância concentrada = explicação simples.
\end{infobox}

A \textbf{entropia} de uma lista de frações $p_1, \dots, p_d$ é $H = -\sum_i p_i \ln p_i$. Ela vale 0 quando toda a importância está em um único atributo e atinge o máximo, $\ln d$, quando todos os atributos têm a mesma importância.

\section{Um exemplo com 4 atributos}
\begin{center}\small
\begin{tabular}{lccccrr}
\toprule
& \textbf{At. 1} & \textbf{At. 2} & \textbf{At. 3} & \textbf{At. 4} & \textbf{Entropia} & \textbf{Nº efetivo} \\
\midrule
Explicação concentrada & 0,97 & 0,01 & 0,01 & 0,01 & 0,17 & 1,2 \\
Explicação A (do fichamento) & 0,5 & 0,3 & 0,1 & 0,1 & 1,17 & 3,2 \\
Explicação uniforme & 0,25 & 0,25 & 0,25 & 0,25 & 1,39 ($=\ln 4$) & 4,0 \\
\bottomrule
\end{tabular}
\end{center}

O valor da entropia é pouco intuitivo. Por isso, neste documento ele também é mostrado como \textbf{número efetivo de atributos}, $e^{H}$: quantos atributos, todos com a mesma importância, dariam a mesma entropia. A explicação concentrada “usa” cerca de 1 atributo; a uniforme usa os 4. Essa conversão é uma ajuda de leitura deste documento e não está no artigo.

\section{O que a métrica não mede}
\begin{itemize}
\item \textbf{Não diz se a explicação está certa.} Uma explicação pode ser bem concentrada no atributo errado. A complexidade fala da \textbf{legibilidade}; a fidelidade (Faithfulness Correlation e Estimate) fala da \textbf{correção}. As duas se complementam.
\item \textbf{Ignora o sinal.} Importância positiva e negativa contam igual, pelo valor absoluto.
\item \textbf{Depende do que se chama de atributo.} Com 4 atributos, o máximo é $\ln 4 = 1{,}39$; com 180\,000 pontos, é $\ln 180\,000 = 12{,}1$. Valores só são comparáveis com o mesmo conjunto de atributos.
\end{itemize}

\partbar{Parte II · O que o artigo diz}

\begin{citacard}{\pagechip{p. 1} §1 Introdução}
\begin{origquote}“Since humans cannot process a lot of information at once, some have argued for minimal model explanations that contain only relevant and representative features [...]; therefore, an explanation should not be complex (i.e., use few features).”\end{origquote}
\begin{tradquote}“Como os humanos não conseguem processar muita informação de uma vez, alguns defendem explicações mínimas, que contenham apenas atributos relevantes e representativos [...]; portanto, uma explicação não deve ser complexa (isto é, deve usar poucos atributos).”\end{tradquote}
\explab{Explicação}
A motivação é a \textbf{compreensão humana}, e não a correção. É um critério sobre quem lê a explicação.
\end{citacard}

\begin{citacard}{\pagechip{p. 2} §3.3 Desideratum: Low Complexity}
\begin{origquote}“A complex explanation is one that uses all d features in its explanation of which features of x are important to f. Though this explanation may be faithful to the model (as defined above), it may be too difficult for the user to understand (especially if d is large).”\end{origquote}
\begin{tradquote}“Uma explicação complexa é aquela que usa todos os d atributos para explicar quais atributos de x são importantes para f. Embora essa explicação possa ser fiel ao modelo (como definido acima), ela pode ser difícil demais de entender para o usuário (especialmente se d for grande).”\end{tradquote}
\explab{Explicação}
O próprio artigo separa as duas coisas: uma explicação pode ser fiel e ainda assim complexa. No ECG, $d = 180\,000$ pontos é exatamente o caso “d grande”.
\end{citacard}

\begin{citacard}{\pagechip{p. 2} §3.3 e Definição 4}
\begin{origquote}“If every feature had equal attribution, the explanation would be complex (even if it is faithful). The simplest explanation would be concentrated on one feature. We define complexity as the entropy of P\textsubscript{g}.”\end{origquote}
\begin{tradquote}“Se todos os atributos tivessem a mesma atribuição, a explicação seria complexa (mesmo que fiel). A explicação mais simples estaria concentrada em um único atributo. Definimos a complexidade como a entropia de P\textsubscript{g}.”\end{tradquote}
\explab{Fórmula}
$P_g(i) = \dfrac{|g(f,x)_i|}{\sum_{j} |g(f,x)_j|}$, \qquad $\mu_C(f, g; x) = -\sum_{i=1}^{d} P_g(i) \ln P_g(i)$.

$P_g(i)$ é a fração da importância total (em valor absoluto) que cabe ao atributo $i$; $\mu_C$ é a entropia dessas frações, calculada para cada exemplo $x$.
\end{citacard}

\begin{citacard}{\pagechip{p. 5} §7 Experimentos}
\begin{origquote}“Before doing aggregation, we unit norm all explanations. For the complexity criterion, we take the positive ℓ1 norm.”\end{origquote}
\begin{tradquote}“Antes da agregação, normalizamos todas as explicações. Para o critério de complexidade, usamos a norma ℓ1 positiva.”\end{tradquote}
\explab{Explicação}
É a divisão pela soma dos valores absolutos que aparece em $P_g$. O artigo não detalha mais nada; em particular, não discute o que seria um atributo em entradas como séries temporais.
\end{citacard}

\begin{citacard}{\pagechip{p. 6} Tabela 1, Tabela 4 e §7.3}
\begin{origquote}“Note that achieving lowest complexity on an image would imply that all attribution is placed on a single pixel.” [Tabela 1] \quad “We show an example with images since it illustrates the notion of lower complexity well; however, other data types (tabular) might be better suited for complexity optimization.” [§7.3]\end{origquote}
\begin{tradquote}“Note que atingir a menor complexidade em uma imagem significaria colocar toda a atribuição em um único pixel.” [Tabela 1] \quad “Mostramos um exemplo com imagens porque ele ilustra bem a noção de complexidade menor; no entanto, outros tipos de dados (tabulares) podem ser mais adequados para otimizar a complexidade.” [§7.3]\end{tradquote}
\explab{Explicação}
Nas imagens do MNIST (784 pixels), a complexidade fica entre 3,2 e 3,7. Na Tabela 4, as explicações SHAP têm complexidade 1,94 (Adult), 1,36 (Iris) e 2,33 (MIMIC), em dados tabulares com poucos atributos. Os próprios autores reconhecem que, com pontos isolados (pixels), a complexidade mínima não faz sentido: concentrar tudo em um pixel não é uma boa explicação. É o mesmo problema dos pontos isolados do ECG.
\end{citacard}

\begin{citacard}{\pagechip{p. 6} §8 Conclusão}
\begin{origquote}“Though we consider one criterion at a time, future work could further axiomatize our criteria, explore the interaction between different evaluation criteria [...].”\end{origquote}
\begin{tradquote}“Embora consideremos um critério por vez, trabalhos futuros poderiam axiomatizar melhor os critérios e explorar a interação entre diferentes critérios de avaliação [...].”\end{tradquote}
\explab{Explicação}
No artigo, a complexidade serve principalmente para \textbf{comparar e combinar funções de explicação} (§5, p.~3), e não para comparar classes de um mesmo modelo, como no TCC. Olhar a complexidade junto com a fidelidade, como se propõe aqui, é justamente a “interação entre critérios” deixada como trabalho futuro.
\end{citacard}

\partbar{Parte III · Aplicação ao TCC}

\section{Do artigo ao código do Quantus}
Arquivo: \texttt{quantus/metrics/complexity/complexity.py} (Quantus 0.6.0). A métrica não está no pipeline atual do TCC (\texttt{src/evaluate.py}).
\begin{center}\small
\begin{tabularx}{\linewidth}{lXX}
\toprule
\textbf{Etapa} & \textbf{Artigo} & \textbf{Quantus} \\
\midrule
Atributo & Cada atributo $x_i$ & Cada valor do mapa achatado (180\,000 no ECG) \\
Frações & $|g_i| / \sum_j |g_j|$ & Igual (\texttt{abs=True}) \\
Entropia & $-\sum P \ln P$ & Igual, com $10^{-12}$ dentro do log \\
Normalização prévia & Não altera o resultado & \texttt{normalise=True} (também não altera: a entropia das frações não depende da escala) \\
\bottomrule
\end{tabularx}
\end{center}
O Quantus é fiel à Definição 4. O problema está na escolha do atributo: com 180\,000 pontos isolados, a métrica herda os mesmos defeitos que levaram a reimplementar a Consistency e a Faithfulness Correlation.

\section{Diagnóstico com os 622 mapas do modelo B}
A complexidade foi calculada nos 622 mapas SHAP do modelo B (predições corretas do teste) de três formas, que diferem só no que se chama de atributo:
\begin{itemize}
\item \textbf{(P) Pontos, como no Quantus:} cada um dos 180\,000 valores do mapa, incluindo o trecho de preenchimento.
\item \textbf{(R) Regiões:} as 36 regiões do batimento (12 derivações × pré-QRS, QRS e ST-T), com a importância de cada uma igual à \textbf{soma do mapa} na região, exatamente como na Consistency, na Faithfulness Correlation e na Faithfulness Estimate.
\item \textbf{(R-abs) Regiões com valores absolutos:} as mesmas 36 regiões, mas somando o valor absoluto de cada ponto.
\end{itemize}
O cálculo usa só os mapas (não chama o modelo) e levou 9 segundos na CPU.

\begin{center}\small
\begin{tabularx}{\linewidth}{Xrrrr}
\toprule
& \textbf{Máximo possível} & \textbf{Média} & \textbf{Nº efetivo médio} & \textbf{Faixa do nº efetivo} \\
\midrule
(P) Pontos (Quantus) & 12,10 & 11,03 & 62\,500 de 180\,000 & 33\,000 a 96\,000 \\
(R) Regiões, soma do mapa & 3,58 & 3,18 & \textbf{24,2 de 36} & \textbf{15,3 a 30,7} \\
(R-abs) Regiões, soma de $|\cdot|$ & 3,58 & 3,53 & 34,1 de 36 & 27,9 a 36 \\
\bottomrule
\end{tabularx}
\end{center}

\subsection{A versão do Quantus mede o comprimento do sinal}\label{sec:pontos}
Os ECGs têm durações diferentes: o trecho real ocupa em média 7\,500 dos 15\,000 instantes, e o resto é preenchimento. Na versão (P), quanto mais longo o ECG, maior a complexidade, porque há mais pontos com alguma importância. A correlação de postos entre complexidade e comprimento é 0,45 com o preenchimento e chega a \textbf{0,88} quando se exclui o preenchimento. Ou seja, a métrica ordena os mapas quase só pela duração do exame. Além disso, em média \textbf{15\% da importância total} do mapa (em valor absoluto) fica no trecho de preenchimento, que não tem informação clínica. A versão (P), por isso, não serve para o TCC.

\subsection{Com regiões, a métrica distingue os mapas}\label{sec:regioes}
Na versão (R), o número de regiões é sempre 36 (todos os ECGs têm os três tipos de região nas 12 derivações), o preenchimento fica de fora e a complexidade praticamente não depende da duração (correlação de postos de 0,07). Os mapas usam, em média, o equivalente a 24 das 36 regiões, com variação real entre mapas (de 15 a 31).

Na versão (R-abs), quase todos os mapas parecem uniformes (34 de 36). Somar valores absolutos acumula o ruído ponto a ponto do GradientSHAP em todas as regiões e esconde as diferenças. A versão (R) é a coerente com as demais métricas do TCC, que também usam a soma do mapa como importância da região. Há um efeito a registrar: dentro de uma região, partes positivas e negativas se cancelam, e uma região com importância “mista” pesa pouco.

\section{Resultados por classe (versão por regiões)}
Classes ordenadas da explicação mais concentrada para a mais espalhada. “Top 3” é a fração da importância total que está nas 3 regiões mais importantes do mapa. A Faithfulness Correlation é a do modelo B (baseline zero, $|S| = 9$).
\begin{center}\small
\begin{tabularx}{\linewidth}{Xrrrrrrr}
\toprule
\textbf{Classe} & \textbf{Mapas} & \textbf{Nº efetivo} & \textbf{Desvio} & \textbf{Top 3} & \textbf{FaithCorr} & \textbf{F1} & \textbf{Posição no F1} \\
\midrule
BRD  & 172 & 22,4 & 2,7 & 33\% & 0,75 & 0,927 & 2.º \\
AF   & 108 & 22,8 & 2,1 & 30\% & 0,54 & 0,931 & 1.º \\
IAVB &  61 & 23,3 & 2,3 & 29\% & 0,65 & 0,884 & 3.º \\
CVP  &  59 & 24,7 & 2,4 & 27\% & 0,61 & 0,837 & 5.º \\
STE  &  14 & 25,5 & 1,7 & 26\% & 0,60 & 0,596 & 9.º \\
CAP  &  45 & 25,7 & 1,8 & 24\% & 0,52 & 0,750 & 8.º \\
BRE  &  20 & 25,8 & 1,7 & 26\% & 0,65 & 0,870 & 4.º \\
STD  &  65 & 26,9 & 1,5 & 23\% & 0,52 & 0,778 & 7.º \\
SNR  &  78 & 26,9 & 1,6 & 23\% & 0,46 & 0,796 & 6.º \\
\bottomrule
\end{tabularx}
\end{center}

\textbf{Leitura descritiva.} As três classes de maior F1 (AF, BRD e IAVB) são as de explicações mais concentradas (22 a 23 regiões efetivas). SNR e STD têm as mais espalhadas (27). As exceções são STE, de menor F1 mas complexidade intermediária, e BRE, de F1 alto mas explicações espalhadas. Dois cuidados: a diferença entre as classes extremas (cerca de 4,5 regiões) é da mesma ordem do desvio dentro de cada classe (1,5 a 2,7); e STE e BRE têm poucos mapas (14 e 20).

\textbf{Uma ressalva de interpretação.} Ser concentrada não é necessariamente melhor no ECG. Algumas classes são, por natureza, visíveis em várias derivações (por exemplo, o ritmo sinusal aparece em todas). Uma explicação espalhada para SNR pode refletir a própria classe, e não uma explicação ruim. Por isso a complexidade deve ser lida junto com a fidelidade, e não sozinha.

\subsection{Relação com a fidelidade}
Mapa a mapa, explicações mais concentradas tendem a ser mais fiéis: a correlação de postos entre a complexidade (R) e a Faithfulness Correlation nos 622 mapas é $-0{,}41$. Isso é usado aqui apenas como diagnóstico, e não como resultado do TCC. Parte dessa relação pode ser mecânica: quando a importância está concentrada, as somas dos subconjuntos sorteados variam mais, e a correlação fica menos sujeita a ruído. Ainda assim, mostra que as duas métricas não medem exatamente a mesma coisa (a relação é moderada), o que é um argumento para usar ambas.

\subsection{Relação com o piso do GradientSHAP}\label{sec:piso}
Esta seção responde a uma pergunta: parte da complexidade medida pode vir do próprio método de explicação, e não do modelo? Para isso, é preciso entender três ideias: os sorteios do GradientSHAP, o piso e o efeito do erro sobre a entropia.

\textbf{(1) O GradientSHAP estima o mapa com 50 sorteios.} O GradientSHAP não calcula o mapa de forma exata; ele \textbf{estima uma média}. Em cada sorteio, ele:
\begin{enumerate}
\item escolhe ao acaso uma \textbf{baseline}, isto é, um ECG de referência, entre as 32 do conjunto de baselines (ver o quadro abaixo);
\item soma ao ECG um pequeno ruído gaussiano (desvio 0,09);
\item escolhe ao acaso um ponto no caminho entre a baseline e o ECG;
\item calcula, nesse ponto, o gradiente da saída da classe e o multiplica por (ECG $-$ baseline).
\end{enumerate}
O mapa salvo é a \textbf{média de 50 desses termos} (\texttt{n\_samples} $= 50$ em \texttt{explain.py}). Cada termo isolado varia bastante, porque depende da baseline e do ponto sorteados. A média de 50 termos é uma boa aproximação, mas não é exata: com outros 50 sorteios (outra \textbf{semente} do gerador aleatório), sai um mapa um pouco diferente. Exemplo com uma única região cujo valor “verdadeiro” (a média com infinitos sorteios) é 0:
\begin{center}\small
\begin{tabular}{lcl}
\toprule
\textbf{Sorteios} & \textbf{Valor da região no mapa} & \textbf{Leitura} \\
\midrule
50 sorteios, semente $s$ & $+0{,}03$ & parece levemente a favor da classe \\
50 sorteios, semente $s+1$ & $-0{,}02$ & parece levemente contra \\
Infinitos sorteios & $0$ & a região não importa \\
\bottomrule
\end{tabular}
\end{center}
Essa diferença não vem do modelo nem do ECG: é o \textbf{erro de estimativa} da média.

\begin{infobox}[note]{As 32 baselines e as sementes}
O GradientSHAP explica o ECG em comparação com ECGs de referência: a pergunta é o que, neste ECG, faz a saída diferir da saída para um ECG “típico”. As \textbf{32 baselines} são 32 ECGs do conjunto de treino (de 4\,812), escolhidos ao acaso e que, juntos, representam esse ECG típico. O número 32 é o padrão de \texttt{src/explain.py} (\texttt{-{}-n-baselines 32}); é uma escolha de configuração do método, entre representatividade e custo, sem justificativa formal registrada. Há dois níveis de sorteio:
\begin{center}\small
\begin{tabularx}{\linewidth}{XXX}
\toprule
\textbf{O que é sorteado} & \textbf{Quando} & \textbf{Semente} \\
\midrule
Quais 32 ECGs do treino formam o conjunto de baselines & Uma única vez & Sempre a mesma (42, a semente do projeto) \\ \addlinespace
Em cada um dos 50 sorteios de um mapa: qual das 32 usar, o ruído e o ponto no caminho & A cada mapa & Uma semente própria de cada mapa \\
\bottomrule
\end{tabularx}
\end{center}
As 32 baselines são, portanto, as mesmas para os 622 mapas e para todas as métricas que refazem o GradientSHAP (sensibilidade, Random Logit). O que muda de um mapa para outro, e entre as duas execuções do piso, são os 50 sorteios feitos dentro desse conjunto fixo.
\end{infobox}

\textbf{(2) O piso mede o tamanho desse erro.} A análise da sensibilidade (Max/Avg-Sensitivity) gerou, para cada um dos 622 mapas, dois mapas do \textbf{mesmo ECG} e da \textbf{mesma classe}, que diferem só nos sorteios (sementes $s$ e $s+1$). A distância entre eles, nas 36 regiões e com os mapas normalizados para norma 1, é o \textbf{piso}. Como nada mudou além dos sorteios, o piso é a variação que o método produz sozinho. A média é 0,355 (a distância vai de 0, mapas iguais, a 2, mapas opostos).

\textbf{Piso de uma classe.} Cada mapa tem o seu piso; o piso de uma classe é a \textbf{média dos pisos dos seus mapas}. Por exemplo, o piso de BRD (0,32) é a média dos 172 pisos dos mapas de BRD.

\textbf{Por que o piso muda de uma classe para outra, se o erro é do GradientSHAP.} O erro vem do método (a média de só 50 termos), mas o \textbf{tamanho} do erro depende de quanto os 50 termos discordam entre si, e isso depende de como o modelo responde para aquela classe. Cada termo é o gradiente num ponto sorteado vezes (ECG $-$ baseline). Se, para aquela classe, o modelo responde de forma parecida qualquer que seja a baseline ou o ponto sorteado, os termos se parecem e a média de 50 já é estável: piso baixo. Se a resposta muda muito conforme o sorteio (por exemplo, o gradiente muda bruscamente ao longo do caminho), os termos discordam e a média de 50 oscila: piso alto. É como estimar uma média com 50 lançamentos: com uma moeda, a estimativa sai estável; com um dado de 20 faces, oscila bem mais. O número de lançamentos é o mesmo; muda o que está sendo sorteado. Assim, o piso é uma propriedade da combinação \textbf{método $\times$ modelo $\times$ classe}: não é um defeito clínico da classe, nem só do GradientSHAP (com mais sorteios ele cairia em todas as classes, mas as diferenças entre elas tenderiam a continuar).

Dois resultados ajudam a ler os pisos por classe (diagnóstico):
\begin{itemize}
\item \textbf{O tamanho da saída explicada não explica o piso.} A soma de um mapa acompanha a diferença entre a saída do modelo para o ECG e a saída média para as baselines. A correlação de postos entre essa soma e o piso, nos 622 mapas, é $-0{,}07$: praticamente nula.
\item \textbf{Dentro de cada classe, o piso varia bastante.} O desvio é de cerca de 0,10 na maioria das classes, enquanto as médias vão de 0,23 (BRE) a 0,46 (CAP). A exceção clara é BRE, com piso baixo e pouca variação (desvio 0,05); nas demais, as diferenças entre classes são modestas frente à variação entre mapas.
\end{itemize}
Por que exatamente uma classe tem piso maior que outra não foi medido; a explicação acima (os termos discordam mais em umas classes do que em outras) é a leitura mais provável, mas não foi verificada.

\textbf{(3) Por que o erro afeta a complexidade.} O erro de estimativa põe pequenas importâncias, positivas ou negativas, em regiões que o modelo talvez nem use. Como a complexidade soma as importâncias em valor absoluto, essas sobras \textbf{espalham} o mapa e aumentam a entropia. Exemplo com 4 regiões:
\begin{center}\small
\begin{tabular}{lrrrrr}
\toprule
& \textbf{Região 1} & \textbf{Região 2} & \textbf{Região 3} & \textbf{Região 4} & \textbf{Nº efetivo} \\
\midrule
Mapa sem erro (infinitos sorteios) & 0,70 & 0,20 & 0,10 & 0,00 & 2,2 \\
Mapa com erro (50 sorteios) & 0,65 & 0,20 & 0,10 & 0,05 & 2,7 \\
\bottomrule
\end{tabular}
\end{center}
O modelo é o mesmo nas duas linhas; só o erro de estimativa fez o mapa parecer usar mais regiões. Portanto, a complexidade medida mistura duas coisas: \textbf{o quanto o modelo espalha a evidência} e \textbf{o ruído do GradientSHAP}. Foi esse mesmo efeito, acumulado ponto a ponto, que deixou a versão (R-abs) quase uniforme (seção~\ref{sec:regioes}).

\textbf{O que o cruzamento mostrou.} Se o ruído fosse a causa principal, os mapas com piso maior seriam sempre os mais espalhados. Para verificar, cruzaram-se os 622 mapas da complexidade com o piso calculado na sensibilidade (mesmos ECGs e classes). Como diagnóstico, e não como resultado do TCC:
\begin{center}\small
\begin{tabular}{lrrrrrrrrr}
\toprule
& \textbf{BRD} & \textbf{AF} & \textbf{IAVB} & \textbf{CVP} & \textbf{STE} & \textbf{CAP} & \textbf{BRE} & \textbf{STD} & \textbf{SNR} \\
\midrule
Nº efetivo & 22,4 & 22,8 & 23,3 & 24,7 & 25,5 & 25,7 & 25,8 & 26,9 & 26,9 \\
Piso & 0,32 & 0,35 & 0,36 & 0,32 & 0,35 & \textbf{0,46} & \textbf{0,23} & 0,40 & 0,41 \\
\bottomrule
\end{tabular}
\end{center}
\begin{itemize}
\item \textbf{Há relação, mas moderada.} Mapa a mapa, a correlação de postos entre o número efetivo de regiões e o piso é 0,38: mapas com mais erro de estimativa tendem a ser um pouco mais espalhados.
\item \textbf{O ruído não explica a ordem das classes.} BRE tem o \textbf{menor} piso de todas (0,23) e mesmo assim explicações espalhadas (25,8). CAP tem o \textbf{maior} piso (0,46) e complexidade intermediária. SNR e STD combinam piso alto e explicações espalhadas, e nelas parte da complexidade pode vir do ruído.
\end{itemize}
\textbf{Como entra na leitura.} O ruído explica uma parte do espalhamento, mas não a ordem das classes. Por isso, a complexidade é usada como está, sem nenhuma correção, e o piso de cada classe é reportado ao lado, como contexto: quando uma classe tem explicações espalhadas \textit{e} piso alto (SNR, STD), parte do espalhamento pode ser do método; quando tem piso baixo (BRE), o espalhamento é mais provavelmente do modelo. A Consistency mistura as mesmas duas causas.

\textbf{Refazer os mapas com mais sorteios.} A única forma de reduzir o erro de estimativa é aumentar o número de sorteios da média (\texttt{n\_samples}). O erro de uma média cai aproximadamente com $1/\sqrt{n}$: com 200 sorteios em vez de 50 (4 vezes mais), o erro cai mais ou menos pela metade. O piso diminuiria, e a complexidade passaria a refletir mais o modelo e menos o método. O custo é regerar os 622 mapas no Kaggle (cerca de 4 vezes o tempo do SHAP atual) e refazer todas as métricas que usam os mapas. Isso foi feito depois, junto com a troca da referência (versão 2 dos mapas, 26/10/2026). As seções 5 a 9 usam os mapas da versão 1, com 50 sorteios; o resultado com 200 sorteios está na seção~\ref{sec:v2}, e a previsão acima não se confirmou.

\subsection{Relação com o eMPRT}\label{sec:emprt}
O eMPRT (Hedström et al., 2023), escolhido na seleção de métricas, mede quanto a entropia da explicação \textbf{sobe} quando os pesos do modelo são todos sorteados de novo: $q = (\xi(\hat{e}) - \xi(e))/\xi(e)$, em que $e$ é o mapa do modelo treinado e $\hat{e}$ o do modelo randomizado. A entropia $\xi$ usada no TCC será a mesma desta métrica, nas 36 regiões. Duas consequências:
\begin{itemize}
\item \textbf{Uma função, duas métricas.} O módulo da complexidade fornece a entropia que o eMPRT usa, e os dois resultados ficam diretamente comparáveis.
\item \textbf{A complexidade limita o eMPRT.} A entropia não passa de $\ln 36 = 3{,}58$. Como os mapas do modelo treinado já têm entropia média 3,18, o aumento $q$ é de no máximo cerca de 13\% (de 5\% a 32\% conforme o mapa). O teto também muda com a classe: cerca de 15\% para BRD, AF e IAVB (mapas mais concentrados) e só 9\% para SNR e STD (os mais espalhados). Sem levar isso em conta, um $q$ menor em SNR poderia parecer pior explicação, quando reflete apenas menos espaço para subir. Por isso, no eMPRT, $q$ será lido junto com o teto de cada mapa.
\end{itemize}

\subsection{Achado lateral: onde está a região mais importante}
O diagnóstico também registrou qual região concentra mais importância em cada mapa. Isso não é complexidade, mas ajuda a interpretar os resultados.
\begin{center}\small
\begin{tabularx}{\linewidth}{lX}
\toprule
\textbf{Classe} & \textbf{Região mais importante (fração dos mapas)} \\
\midrule
BRD  & V1 QRS (87\%) \\
BRE  & V2 QRS (50\%), V1 QRS (25\%) \\
IAVB & aVR pré-QRS (62\%), II pré-QRS (10\%), V1 pré-QRS (7\%) \\
STD  & aVR ST-T (37\%), II ST-T (9\%), V2 ST-T (9\%) \\
SNR  & aVR pré-QRS (47\%), V1 QRS (18\%) \\
CVP  & aVR QRS (29\%), aVF QRS (17\%), V2 QRS (10\%) \\
STE  & aVR pré-QRS (57\%), V1 QRS (21\%) \\
CAP  & V2 QRS (13\%), V1 pré-QRS (13\%), aVR pré-QRS (11\%) \\
AF   & V5 QRS (17\%), V4 QRS (16\%), aVR QRS (14\%) \\
\bottomrule
\end{tabularx}
\end{center}
Vários casos batem com os critérios clássicos de diagnóstico: o bloqueio de ramo direito (BRD) se vê no QRS de V1; o de ramo esquerdo (BRE), no QRS de V1 e V2; o bloqueio AV de 1.º grau (IAVB) é um intervalo PR longo, que fica na região pré-QRS; o infradesnivelamento (STD), no segmento ST. Outros chamam atenção, como STE (supradesnivelamento de ST), cuja região principal é pré-QRS e não ST-T, e a classe de menor F1. Essa leitura clínica precisa ser confirmada com a literatura e com o orientador. Uma análise formal de \textbf{plausibilidade} (localização com máscaras clínicas) foi avaliada na seleção de métricas e deixada de fora: ela mede se a explicação coincide com o critério clínico, e não se é fiel ao modelo, e só daria para definir máscaras para cerca de 5 das 9 classes. A tabela fica apenas como apoio à interpretação.

\section{Qual métrica de complexidade usar}\label{sec:escolha}
O Quantus traz três métricas de complexidade. As três medem a mesma propriedade (quão concentrada está a importância) e diferem só na conta. As definições abaixo vêm do código do Quantus 0.6.0 (\texttt{quantus/metrics/complexity/}); os artigos de Chalasani et al.\ (2020) e de Nguyen e Martínez (2020) não foram lidos.
\begin{center}\small
\begin{tabularx}{\linewidth}{lXXX}
\toprule
& \textbf{Complexity} (Bhatt et al., 2020) & \textbf{Sparseness} (Chalasani et al., 2020) & \textbf{Effective Complexity} (Nguyen e Martínez, 2020) \\
\midrule
Conta & Entropia das frações de $|$importância$|$ & Índice de Gini dos valores de $|$importância$|$ & Quantos valores de $|$importância$|$ passam de um limiar $\varepsilon$ \\ \addlinespace
Escala & 0 a $\ln 36$; vira nº efetivo de regiões (1 a 36) & 0 (uniforme) a 1 (tudo em uma região) & 0 a 36 regiões \\ \addlinespace
Depende da escala do mapa? & Não (usa frações) & Não & \textbf{Sim} \\ \addlinespace
Parâmetro a escolher & Nenhum & Nenhum & \textbf{O limiar} $\varepsilon$ \\ \addlinespace
No Quantus, para o ECG & Pontos com preenchimento (mede a duração) & Pontos com preenchimento & Pontos com preenchimento; $\varepsilon = 10^{-5}$ após dividir pelo máximo, o que conta praticamente todo ponto não nulo \\
\bottomrule
\end{tabularx}
\end{center}
Qualquer uma delas precisaria da versão por regiões. A escolha entre elas:
\begin{itemize}
\item \textbf{Complexity (escolhida).} Não tem parâmetro, não depende da escala e tem leitura intuitiva pelo número efetivo de regiões (“a explicação de BRD usa o equivalente a 22 regiões; a de SNR, 27”). É também a entropia do eMPRT (seção~\ref{sec:emprt}), e o diagnóstico já foi feito com ela.
\item \textbf{Sparseness (verificação).} O índice de Gini também não tem parâmetro nem depende da escala. A diferença é de ênfase: o Gini é mais sensível a quanto as primeiras regiões concentram; a entropia considera a distribuição inteira. O Gini foi calculado no mesmo módulo, como verificação, e \textbf{confirma a entropia}: a ordem das classes é a mesma, exceto por uma troca entre BRD e AF no topo, com diferença desprezível; mapa a mapa, a correlação de postos entre as duas é $-0{,}99$ (seção~\ref{sec:impl}). A escolha entre entropia e Gini, portanto, não muda o resultado.
\item \textbf{Effective Complexity (fora).} Depende de um limiar sobre o valor da importância. A soma de um mapa do GradientSHAP acompanha a diferença de saída do modelo entre o ECG e as baselines, que varia entre ECGs e classes; com um limiar fixo, classes de saída maior teriam mais regiões acima dele mesmo com a mesma concentração. Dividir pelo máximo, como no Quantus, troca esse problema pela escolha de um limiar relativo arbitrário. E com 36 regiões o resultado é uma contagem grosseira, com poucos valores possíveis.
\end{itemize}

\section{Decisão}\label{sec:decisao}
\begin{center}\small
\begin{tabularx}{\linewidth}{XX}
\toprule
\textbf{A favor} & \textbf{Cuidados} \\
\midrule
Vem do mesmo artigo da Faithfulness Correlation e da Max/Avg-Sensitivity, e completa os três critérios propostos. & Não mede correção, e sim legibilidade. Menor não é automaticamente melhor no ECG: algumas classes são naturalmente espalhadas. \\ \addlinespace
Custo quase nulo: 9 segundos para os 622 mapas, sem chamar o modelo. & A versão do Quantus (pontos) não serve: mede o comprimento do sinal (0,88) e inclui o preenchimento. \\ \addlinespace
Sem baseline e sem perturbação: não sofre dos artefatos que afetaram as versões antigas da Faithfulness Correlation e da Estimate. & O resultado depende da divisão em 36 regiões e do uso da soma do mapa (com cancelamento entre positivos e negativos), escolhas que precisam ser justificadas. \\ \addlinespace
Leitura intuitiva por regiões: “a explicação usa o equivalente a 24 das 36 regiões”. & Diferenças entre classes modestas frente à variação dentro de cada classe. \\ \addlinespace
Mede algo diferente da fidelidade (relação moderada, $-0{,}41$). & É de outra categoria do Quantus (complexidade); incluí-la amplia o escopo do TCC. \\ \addlinespace
Fornece a entropia do eMPRT e o teto de cada mapa para lê-lo (seção~\ref{sec:emprt}). & Parte da complexidade poderia vir do erro de estimativa do GradientSHAP (relação de 0,38 com o piso; seção~\ref{sec:piso}). Com 200 sorteios, o erro caiu pela metade e a complexidade não mudou (seção~\ref{sec:v2}). \\ \addlinespace
\bottomrule
\end{tabularx}
\end{center}

\begin{infobox}[ok]{Decisão (22/10/2026)}
\textbf{A complexidade entra no TCC como métrica complementar, na versão por regiões (R)}, junto com Random Logit e eMPRT (documento \textit{Seleção de métricas}). O objetivo não é verificar se a explicação é clinicamente compatível, mas comparar o quanto as explicações de cada classe são concentradas ou espalhadas, o que dá base para discutir os resultados de fidelidade, robustez e randomização. A métrica é sempre lida ao lado da fidelidade (uma boa explicação é fiel \textit{e} concentrada) e do piso. A \texttt{quantus.Complexity} ponto a ponto não é usada. A ampliação do escopo será apresentada ao orientador junto com a mudança de rumo da seleção.
\end{infobox}

\section{Implementação e resultados}\label{sec:impl}
Módulo \texttt{src/complexity.py}, no mesmo padrão de \texttt{src/consistency.py} e \texttt{src/sensitivity.py}:
\begin{itemize}
\item \textbf{Entrada:} os 622 mapas do modelo B, já salvos. Não chama o modelo, roda na CPU em cerca de 10 segundos e não precisa do Kaggle.
\item \textbf{Regiões:} as mesmas 36 da Faithfulness Correlation (função \texttt{region\_index}), com a importância da região igual à soma do mapa.
\item \textbf{Por mapa:} entropia $H$, número efetivo de regiões $e^{H}$, fração nas 3 regiões mais importantes, região mais importante e, como verificação, o índice de Gini.
\item \textbf{Por classe:} média e desvio do número efetivo de regiões e do Gini, com o piso da sensibilidade ao lado como contexto e a ordem comparada à do F1 de forma descritiva (sem correlação com o F1).
\item \textbf{Saída:} \texttt{results/modeloB\_complexity/} (tabela por mapa e por classe), seguindo o padrão das outras métricas.
\end{itemize}
A função da entropia por regiões fica exposta para ser reaproveitada pelo eMPRT.

\textbf{Execução (23/10/2026).} Os valores coincidem com os do diagnóstico da seção~5 (diferença máxima de $5 \times 10^{-6}$ na entropia, por arredondamento, e a mesma região mais importante em todos os mapas). Classes da explicação mais concentrada para a mais espalhada:
\begin{center}\small
\begin{tabular}{lrrrrrr}
\toprule
\textbf{Classe} & \textbf{Mapas} & \textbf{Entropia} & \textbf{Nº efetivo (de 36)} & \textbf{Gini} & \textbf{Top 3} & \textbf{Piso} \\
\midrule
BRD  & 172 & 3,10 & 22,4 & 0,508 & 32,7\% & 0,32 \\
AF   & 108 & 3,12 & 22,8 & 0,514 & 29,9\% & 0,35 \\
IAVB &  61 & 3,14 & 23,3 & 0,503 & 29,3\% & 0,36 \\
CVP  &  59 & 3,20 & 24,7 & 0,469 & 27,1\% & 0,32 \\
STE  &  14 & 3,24 & 25,5 & 0,450 & 25,8\% & 0,35 \\
CAP  &  45 & 3,25 & 25,7 & 0,445 & 24,4\% & 0,46 \\
BRE  &  20 & 3,25 & 25,8 & 0,443 & 26,4\% & 0,23 \\
STD  &  65 & 3,29 & 26,9 & 0,414 & 22,9\% & 0,40 \\
SNR  &  78 & 3,29 & 26,9 & 0,414 & 23,0\% & 0,41 \\
\midrule
Global & 622 & 3,18 & 24,2 & 0,475 & 28,1\% & 0,36 \\
\bottomrule
\end{tabular}
\end{center}
\begin{itemize}
\item \textbf{Gini como verificação.} Gini maior significa explicação mais concentrada, o oposto da entropia. A ordem pelo Gini é a mesma da entropia, exceto pela troca entre BRD (0,508) e AF (0,514). Mapa a mapa, a correlação de postos entre as duas medidas é $-0{,}99$. O resultado não depende de qual das duas é usada.
\item \textbf{Leitura frente ao F1 (descritiva).} Valem as observações da seção~6: as três classes de maior F1 (AF, BRD e IAVB) têm as explicações mais concentradas; SNR e STD, as mais espalhadas; STE e BRE são as exceções. A diferença entre os extremos (4,5 regiões) é da ordem do desvio dentro das classes, e o piso (seção~\ref{sec:piso}) mostra que parte do espalhamento pode vir do erro de estimativa do GradientSHAP.
\end{itemize}

\section{Resultados com os mapas da versão 2}\label{sec:v2}
\textbf{O que mudou nos mapas.} Desde 26/10/2026 o TCC usa a versão 2 dos mapas SHAP, a opção D do documento \url{fichamentos/verificacao-referencia-gradientshap}. A referência (o conjunto de baselines) deixou de ser 32 ECGs sorteados do treino e passou a ser \textbf{balanceada}: 4 ECGs de rótulo único de cada uma das 9 classes, 36 no total. Cada mapa passou a ser a média de \textbf{200 sorteios} (4 blocos de 50), em vez de 50. Com isso, as seções anteriores continuam valendo como explicação do método, mas os números atuais são os desta seção. A complexidade foi recalculada na CPU com os mapas novos (\texttt{results/modeloB\_v2/complexity/}), e o piso de cada mapa vem da sensibilidade refeita com a mesma configuração (Kaggle, versão 24 do runner, 27/10/2026).

\begin{center}\small
\begin{tabularx}{\linewidth}{Xrrrrrrr}
\toprule
\textbf{Classe} & \textbf{Mapas} & \textbf{Nº efetivo} & \textbf{Desvio} & \textbf{Gini} & \textbf{Top 3} & \textbf{Piso} & \textbf{Nº efetivo, versão 1} \\
\midrule
BRD  & 172 & 22,2 & 2,3 & 0,514 & 32,9\% & 0,13 & 22,4 \\
AF   & 108 & 22,9 & 2,3 & 0,513 & 29,6\% & 0,19 & 22,8 \\
IAVB &  61 & 23,0 & 2,2 & 0,508 & 29,7\% & 0,18 & 23,3 \\
CVP  &  59 & 25,0 & 2,4 & 0,460 & 27,1\% & 0,17 & 24,7 \\
CAP  &  45 & 25,4 & 1,9 & 0,451 & 25,8\% & \textbf{0,25} & 25,7 \\
STE  &  14 & 25,5 & 1,7 & 0,453 & 24,8\% & 0,18 & 25,5 \\
SNR  &  78 & 26,2 & 1,3 & 0,433 & 24,1\% & 0,23 & 26,9 \\
BRE  &  20 & 26,2 & 1,8 & 0,435 & 24,6\% & \textbf{0,13} & 25,8 \\
STD  &  65 & 26,9 & 1,7 & 0,416 & 23,0\% & 0,23 & 26,9 \\
\midrule
Global & 622 & 24,1 & 2,7 & 0,479 & 28,4\% & 0,18 & 24,2 \\
\bottomrule
\end{tabularx}
\end{center}
\textit{Classes da explicação mais concentrada para a mais espalhada. Nº efetivo de regiões de 36; piso = distância entre dois mapas do mesmo ECG e classe que diferem só nos sorteios (seção~\ref{sec:piso}).}

\textbf{O que se observa.}
\begin{itemize}
\item \textbf{A complexidade quase não mudou.} A média global foi de 24,2 para 24,1 regiões efetivas, e nenhuma classe mudou mais de 0,8. As três classes mais concentradas continuam BRD, AF e IAVB, e STD continua a mais espalhada. As trocas ficam no meio e no fim da ordem, entre classes separadas por menos de 1 região: SNR (26,9 → 26,2) e BRE (25,8 → 26,2) passam a empatar, e CAP passa à frente de STE. Mapa a mapa, a ordem dos valores das duas versões tem correlação de postos de 0,76, com diferença média de 1,5 região (diagnóstico).
\item \textbf{O piso caiu pela metade} (0,36 → 0,18), como previsto pelo erro de uma média, que cai com $1/\sqrt{n}$. A queda ocorreu em todas as classes, e as diferenças entre elas continuam: CAP tem o maior piso (0,25), e BRE e BRD, os menores (0,13).
\item \textbf{A previsão do fim da seção~\ref{sec:piso} não se confirmou.} Esperava-se que, com menos erro de estimativa, os mapas ficassem mais concentrados. Mas o erro caiu pela metade e a complexidade ficou praticamente igual. No teste da referência (\url{results/modeloB_basebal}), usar 200 sorteios com a referência antiga também deixou a complexidade igual. Logo, o espalhamento dos mapas vem do modelo, e não do ruído do GradientSHAP. A relação mapa a mapa entre complexidade e piso continua moderada (correlação de postos de 0,44; 0,38 na versão 1), mas ela indica que os mapas mais difíceis de estimar são também os mais espalhados, e não que o ruído cause o espalhamento.
\item \textbf{O ruído continua não explicando a ordem das classes.} BRE tem o menor piso (0,13) e explicações das mais espalhadas (26,2); CAP tem o maior piso (0,25) e complexidade intermediária.
\item \textbf{Gini e entropia continuam concordando} (correlação de postos de $-0{,}99$ mapa a mapa), e a relação com a fidelidade se mantém: correlação de postos de $-0{,}38$ entre o número efetivo e a Faithfulness Correlation da versão 2 ($-0{,}41$ na versão 1). Ambas são usadas só como diagnóstico.
\item \textbf{Teto do eMPRT.} A entropia média continua em 3,18, e o aumento máximo possível $q$ continua em cerca de 13\% (de 4\% a 29\% conforme o mapa): cerca de 15\% para BRD, AF e IAVB e de 9\% a 10\% para SNR, STD e BRE.
\end{itemize}

\textbf{Região mais importante (versão 2).} A tabela da seção anterior muda pouco: em 71\% dos mapas, a região mais importante é a mesma nas duas versões. Os padrões ficam, em geral, mais nítidos: BRD em V1 QRS (90\%), BRE em V2 QRS (65\%) e V1 QRS (15\%), IAVB em aVR pré-QRS (69\%) e II pré-QRS (13\%), SNR em aVR pré-QRS (60\%) e STE em aVR pré-QRS (64\%). STD fica dividida entre aVR ST-T (23\%) e aVR pré-QRS (22\%). AF e CAP continuam sem região dominante: aVR pré-QRS aparece em 21\% e 33\% dos mapas.

\textbf{Leitura frente ao F1 (descritiva).} Continua igual à da versão 1: as três classes de maior F1 (AF, BRD e IAVB) têm as explicações mais concentradas; SNR e STD, das mais espalhadas; STE (menor F1, complexidade intermediária) e BRE (F1 alto, explicações espalhadas) são as exceções.

\vspace{0.5em}
{\small \textbf{Resultados} (versionados): versão 2 em \url{results/modeloB_v2/complexity/}; \\ versão 1 em \url{results/modeloB_complexity/} (\texttt{complexity.csv}, \texttt{complexity\_mapas.csv} e \texttt{README.md} com o comando). \\
\textbf{Arquivos do diagnóstico} (pasta não versionada): \\ \texttt{output/leituras/complexity\_diag.py} \\ \texttt{output/leituras/complexity\_diag.csv} (uma linha por mapa)}

\end{document}
